% =========================================================================
% Capítulo 2 — Contexto (OBLIGATORIO)
% Situación en la que nace el proyecto: el problema de las redes sociales
% de entretenimiento vs. una red orientada a cumplir objetivos, el marco
% académico (TFG) y personal, etc.
% (La relación con los ODS va en el capítulo 4, junto con los beneficios,
% según la guía definitiva de la EIB.)
% =========================================================================

\chapter{Contexto}
\label{ch:contexto}

\section{El problema: metas que se abandonan y redes que no acompañan}
\label{sec:contexto-problema}

Las metas vitales se abandonan siguiendo un patrón reconocible: deseos formulados sin plazo ni medida, progreso que nadie registra y un camino que se recorre en solitario, sin nadie alrededor que acompañe o pida cuentas.
El capítulo \ref{ch:beneficios} fundamenta esa lectura con la literatura de la psicología de la motivación; este capítulo la da por presentada para situar el producto.

Las redes sociales de entretenimiento no alivian el problema y hay razones fundadas —presentadas en el capítulo \ref{ch:beneficios}— para pensar que lo agravan: su experiencia gira en torno al desplazamiento infinito por contenido ajeno y a la comparación con vidas editadas.
El hueco que \bukit ocupa es precisamente ese: entre los productos que el capítulo \ref{ch:estado-del-arte} analiza no hay una red social generalista cuya experiencia completa esté orientada a que la persona \emph{haga}, en lugar de a que mire.

\section{La propuesta: una red social orientada a la acción}
\label{sec:contexto-propuesta}

\bukit se organiza en cuatro pestañas, el patrón de navegación asentado por las redes dominantes \parencite{instagram-app}, para que la barrera de entrada sea la de una aplicación ya conocida:

\begin{enumerate}
  \item \textbf{Inicio}: un \eng{feed} —el flujo de publicaciones de la pantalla de inicio— personalizado en mosaico con ideas de metas, progresos, metas completadas y no logradas y vídeos de motivación, más un buscador de perfiles, publicaciones y listas; deslizando, un segundo \eng{feed} solo con el contenido de las personas seguidas.
  \item \textbf{Listas}: el centro de control de los objetivos: crear, ver, editar y eliminar listas de metas con plazo, con recordatorios que envía la propia aplicación.
  \item \textbf{Mensajes y mapa}: la mensajería, con un mini-mapa en la parte superior; el mapa a pantalla completa ofrece dos vistas —la de seguidos mutuos y la de descubrimiento de personas afines, filtrable por país, por radio de cercanía o en todo el mundo— y aquí viven también las notificaciones.
  \item \textbf{Perfil}: quién eres: hasta cinco publicaciones de tema libre —las únicas no ligadas a listas— e historias ilimitadas.
\end{enumerate}

Sobre esa estructura, cinco reglas de producto concretan los tres pilares del concepto —estructura, comunidad y recompensa al avance— y condicionan el diseño técnico del resto de la memoria:

\begin{itemize}
  \item \textbf{Las listas son la unidad de las metas.}
        Una lista agrupa metas con plazo, y las publicaciones de progreso o de meta completada solo pueden crearse desde listas propias: nadie puede presumir de un avance que no persigue de verdad.
        El plan gratuito admite un máximo de tres listas.
  \item \textbf{La visibilidad premia el avance.}
        Hay dos tipos de usuario: normal (cualquier persona registrada) y TOP (creador admitido por formulario externo, con vídeo ilimitado y mayor peso en el \eng{feed}); la ordenación del \eng{feed} pondera el progreso y la aportación, no la mera popularidad, y los vídeos de motivación son exclusivos de los perfiles TOP.
  \item \textbf{El contenido de perfil es efímero.}
        Todo lo publicado en el perfil caduca a los 60 días salvo que su autor pulse «mantener»: el escaparate personal no es el centro de la experiencia.
  \item \textbf{Conocer gente exige verificación.}
        La verificación de identidad —documento de identidad, teléfono y fotografías de validación, delegada en un proveedor especializado— desbloquea publicar desde el perfil y el mapa de descubrimiento, reservado a personas mayores de edad, con ubicación siempre aproximada y con consentimiento explícito.
  \item \textbf{Los recordatorios empujan.}
        La aplicación recuerda los plazos de las metas: el compromiso es medible y con fecha, no un deseo vago.
\end{itemize}

La traducción de estas reglas a beneficios esperados, con su fundamento en la literatura, está en el capítulo \ref{ch:beneficios}; su traducción a requisitos del sistema, en el capítulo \ref{ch:estado-del-arte}.

\section{Marco académico y personal}
\label{sec:contexto-marco}

Este proyecto es el \ac{TFG} del Grado en Ingeniería Informática de Gestión y Sistemas de Información de la \ac{EIB}, centro de la \ac{UPVEHU}, con 12 créditos \ac{ECTS} —unas 360 horas de referencia académica, cuya conciliación con la dedicación realmente planificada se recoge en el capítulo \ref{ch:gantt-tfg}— y desarrollo previsto a lo largo del curso 2026--2027 completo.
El calendario está condicionado por circunstancias personales declaradas en la planificación: compromisos laborales durante el verano, una estancia Erasmus en el primer cuatrimestre y prácticas de empresa desde febrero (capítulos \ref{ch:tareas} y \ref{ch:gantt-tfg}).

El trabajo no parte de cero.
Su precedente directo es TrackBites \parencite{deleanu-trackbites}, la aplicación web completa que el autor construyó con anterioridad y cuya arquitectura funciona como línea base de las decisiones tecnológicas de este proyecto (capítulo \ref{ch:alternativas}).
El desarrollo se realiza en solitario y con la asistencia de \ac{IA} generativa ya acotada en la metodología (apartado \ref{sec:intro-metodologia}).
